% Shared preamble for the MAGI user manual.
\usepackage[T1]{fontenc}
\usepackage[utf8]{inputenc}
\usepackage{lmodern}
\usepackage[a4paper,margin=2.4cm,headheight=15pt]{geometry}
\usepackage{amsmath,amssymb}
\usepackage{graphicx}
\usepackage{booktabs}
\usepackage{longtable}
\usepackage{array}
% [table] pulls in colortbl, needed for the \rowcolor highlighting in the
% per-layer usage table (Visualization section).
\usepackage[table]{xcolor}
\usepackage{listings}
\usepackage{enumitem}
\usepackage{microtype}
\usepackage{caption}
\usepackage{fancyhdr}
\usepackage{titlesec}
\usepackage{tcolorbox}
\usepackage[hidelinks,breaklinks]{hyperref}

\tcbuselibrary{skins,breakable}

% ---------------------------------------------------------------- colours
\definecolor{magiblue}{HTML}{1F4E79}
\definecolor{magigrey}{HTML}{4A4A4A}
\definecolor{codebg}{HTML}{F6F8FA}
\definecolor{codekw}{HTML}{A626A4}
\definecolor{codestr}{HTML}{50A14F}
\definecolor{codecom}{HTML}{9CA0A4}
\definecolor{warnbg}{HTML}{FFF4E5}
\definecolor{warnrule}{HTML}{C77700}
\definecolor{notebg}{HTML}{EEF4FB}

\hypersetup{
  colorlinks=true,
  linkcolor=magiblue,
  urlcolor=magiblue,
  citecolor=magiblue,
}

% ---------------------------------------------------------------- headers
\pagestyle{fancy}
\fancyhf{}
\fancyhead[L]{\small\textcolor{magigrey}{MAGI \magiversion{} --- User Manual}}
\fancyhead[R]{\small\textcolor{magigrey}{\leftmark}}
\fancyfoot[C]{\small\thepage}
\renewcommand{\headrulewidth}{0.4pt}

\titleformat{\section}{\Large\bfseries\color{magiblue}}{\thesection}{0.7em}{}
\titleformat{\subsection}{\large\bfseries\color{magiblue}}{\thesubsection}{0.7em}{}

% ---------------------------------------------------------------- listings
\lstdefinestyle{magi}{
  backgroundcolor=\color{codebg},
  basicstyle=\ttfamily\footnotesize,
  keywordstyle=\color{codekw}\bfseries,
  stringstyle=\color{codestr},
  commentstyle=\color{codecom}\itshape,
  breaklines=true,
  breakatwhitespace=false,
  showstringspaces=false,
  frame=single,
  rulecolor=\color{codebg!60!black!20},
  framesep=5pt,
  xleftmargin=6pt,
  columns=fullflexible,
  keepspaces=true,
  upquote=true,
  literate={-}{{-}}1 {_}{{\_}}1,
}
\lstset{style=magi}
\lstdefinelanguage{shell}{
  keywords={python,pytest,pip,cd,export,bash},
  sensitive=true,
  comment=[l]{\#},
}
\newcommand{\py}[1]{\lstinline[language=Python]|#1|}
\newcommand{\code}[1]{\texttt{\small #1}}
% Long snake_case identifiers inside \code give TeX no break points, which is
% what runs a line into the margin. \cbrk marks a legal break inside one
% (no hyphen inserted); a little \emergencystretch absorbs the rest.
\newcommand{\cbrk}{\discretionary{}{}{}}
\setlength{\emergencystretch}{1.2em}

% ---------------------------------------------------------------- callouts
\newtcolorbox{magiwarn}[1][]{
  breakable, enhanced, colback=warnbg, colframe=warnrule,
  boxrule=0pt, leftrule=3pt, arc=1pt,
  fonttitle=\bfseries, coltitle=warnrule, title={#1},
  left=7pt, right=7pt, top=5pt, bottom=5pt,
}
\newtcolorbox{maginote}[1][]{
  breakable, enhanced, colback=notebg, colframe=magiblue,
  boxrule=0pt, leftrule=3pt, arc=1pt,
  fonttitle=\bfseries, coltitle=magiblue, title={#1},
  left=7pt, right=7pt, top=5pt, bottom=5pt,
}
